\chapter{Importar imagens externas}

\eyebrow{importação de imagens externas}

\vspace{4pt}
Fotos já existentes, tiradas com outra câmera ou em outra sessão, são trazidas
pelo objeto \textbf{Import images} sem quebrar o pareamento.

\section{Duas maneiras de importar}

Uma pasta é selecionada; o aplicativo a percorre, calcula um \tele{sha256} para cada
imagem (de modo que duplicatas sejam detectadas) e lê o EXIF. Em seguida, há dois
modos:

\begin{enumerate}[leftmargin=1.4em,itemsep=4pt]
  \item \textbf{Importar todas como novas carcaças.} Cada foto se torna
        uma nova carcaça, com a etiqueta física extraída do nome do arquivo. Os dados
        reais são preenchidos depois na visão Overview da carcaça. Nomes de arquivo
        repetidos recebem um sufixo numérico automaticamente, de modo que as etiquetas
        permanecem únicas.
  \item \textbf{Parear a uma carcaça existente.} Item a item, cada foto é conciliada a
        uma carcaça já presente no lote, o que é útil quando várias fotos pertencem ao
        mesmo animal.
\end{enumerate}

\begin{note}[title={O pareamento é preservado em ambos os modos}]
Qualquer que seja o modo escolhido, a carcaça é criada (ou escolhida) no momento da
importação, de modo que a imagem não fica órfã. Duplicatas já presentes no conjunto de
dados são ignoradas por hash.
\end{note}

O antigo modo de falha (imagens órfãs às quais depois se atribuíam vínculos
inventados) passa a ser uma etapa explícita, resolvida pelo operador no momento da
importação.
